
% JuliaCon proceedings template
\documentclass{juliacon}
\usepackage{amsmath}

\setcounter{page}{1}
\newcommand{\drv}{\textrm{d}}

\begin{document}

% **************GENERATED FILE, DO NOT EDIT**************

\title{Interplay between Chaos and Stochasticity in Celestial Mechanics}

\author[1]{Matteo Manzi}
\author[2]{Stefano Guidolotti}
\author[2]{Mattia Petrini}
\affil[1]{CrunchDAO}
\affil[2]{University of Pisa}

\keywords{Dynamical Systems, Chaos Theory, Finite-time Lyapunov Exponent, Stochastic differential equation, Lagrangian coherent structures}

\hypersetup{
pdftitle = {Interplay between Chaos and Stochasticity in Celestial Mechanics},
pdfsubject = {JuliaCon 2019 Proceedings},
pdfauthor = {Matteo Manzi, Stefano Guidolotti, Mattia Petrini},
pdfkeywords = {Dynamical Systems, Chaos Theory, Finite-time Lyapunov Exponent, Stochastic differential equation, Lagrangian coherent structures},
}



\maketitle

\begin{abstract}

This work is focused on the development of an open-source Julia-based repository for the anlysis of chaos in dynamical systems, 
in particular for systems described by ordinary and stochastic differential equations, using Finite-Time Lyapunov exponents (FTLE).
The novel application of this scalar field for stochastic processes allows one to generalize the definition of chaos in a probabilistic sense.
This probabilistic generalization is useful for both of uncertainty quantification, and robust trajectory design.
Bifurcating phenomena and invariant sets in time-dependant dynamical systems are discussed, particularly in the context of Lagrangian coherent structures.

\end{abstract}

\section{Introduction}

Chaotic behavior is omnipresent in celestial mechanics dynamical systems and it is relevant for both the understanding and 
leveraging the stability of planetary systems, inner solar system in particular \cite{celletti}. 
The quantification of the probability of impacts of near Earth objects after close encounters with celestial bodies; 
the possibility of designing robust low energy transfer trajectories, not limited to invariant manifolds but also leveraging the 
weak stability boundary for the design of the ballistic captures trajectories in time-dependent dynamical systems; 
the characterization of diffusion processes in Nearly-Integrable Hamiltonian systems in celestial mechanics. 

Building on \cite{Szebehely82}, \cite{ManziAAS2020, VasileManzi} investigated a probabilistic generalization of chaos, in which uncertainty in initial conditions and dynamical parameters is considered.

Here we generalize, inspired by \cite{Balibrea-Iniesta}, in order to introduce a quantification of chaos for stochastic processes.

\section{Metodology}

%slide1 Finite-Time Lyapunov Exponent
\begin{equation}
    \dot{\mathbf{x}}=\mathbf{f}(\mathbf{x}, \mathbf{p})
\end{equation}

\begin{equation}
    \frac{\drv}{\drv t} \frac{\partial \mathbf{x} }{\partial \mathbf{x_0}} = \frac{\partial \mathbf{f}}{\partial \mathbf{x}}\frac{\partial \mathbf{x}}{\partial \mathbf{x_0}}
\end{equation}

\begin{equation}
    \Delta(\mathbf{x}, \mathbf{p}) = \left( \frac{\partial \mathbf{x}}{\partial \mathbf{x_0}}\right)^T \frac{\partial \mathbf{x}}{\partial \mathbf{x_0}}
\end{equation}

\begin{equation}
    \sigma(\mathbf{x}, \mathbf{p}, T)=\frac{1}{|T|}ln(\sqrt{\lambda_{max}(\Delta(\mathbf{x}, \mathbf{p}))})
\end{equation}

%slide2 Standard Map
% One of the most popular mappings in the theory of Dynamical Systems is the so-called \textbf{standard map}, which was
Discrete Dynamical System introduced by Chirikov \cite{CHIRIKOV1979263}, it is defined by the equations:
\begin{equation}
    \begin{cases}
    y'=y+\epsilon f(x) \\
    x'=x+y'
\end{cases}
\end{equation}
where $y\in\mathbb{R}$, $x\in\mathbb{T}\equiv \mathbb{R}/(2\pi\mathbb{Z})$, $\epsilon$ is a positive real parameter, called the perturbing parameter, and $f = f(x)$ is an analytic, periodic function.
\vspace{0.5cm}

%slide4 Perturbed Pendulum
The motion of a periodically-perturbed pendulum is given by:
\begin{equation}
\ddot{x}=(\alpha \cos5t -1 )\sin x
\end{equation}

How is Chaotic Behaviour Influenced by the initial condition of the system? 

And by its parameters $\alpha$ and T?

%slide5 The Double Gyre
The double gyre model consists of a pair of counter-rotating gyres, with a time-periodic perturbation. \\ The system is modelled as a first order differential equation:

\begin{equation}
\dot{\mathbf{z}} =\frac{\drv}{\drv t}\begin{bmatrix}
    x \\
    y  \end{bmatrix}=\pi A\begin{bmatrix}
    -\sin(\pi f(x,t))\cos(\pi y) \\
    \cos(\pi f(x,t))\sin(\pi y)\frac{df}{dx}  \end{bmatrix} 
\end{equation}

%slide6 Stochastic Differential Equations (SDEs)
A stochastic differential equation \cite{oksendal} is known as:
\begin{equation}
    \drv \mathbf{X}=\mathbf{f}(\mathbf{X},t) \drv t + \mathbf{G}(\mathbf{X},t) \drv \mathbf{W_t}
\end{equation}
\begin{equation*}
   \mathbf{X},\mathbf{f} \in \mathbb{R}^{n} 
\end{equation*}
\begin{equation*}
    \mathbf{W_t} \in \mathbb{R}^{m}
\end{equation*}
\begin{equation*}
    \mathbf{G} \in \mathbb{R}^{n \times m}
\end{equation*}

The equation consists of two major components. The deterministic drift denoted as $\mathbf{f}$ and the stochastic diffusion $\mathbf{G}$. $\mathbf{G}$ depends on $m$ Brownian processes expressed as $\mathbf{W_t}$.

%slide7 Finite-Time Lyapunov Exponent for SDEs
\begin{equation}
    \frac{\drv}{\drv t} \frac{\partial \mathbf{X} }{\partial \mathbf{X_0}} = \frac{\partial \mathbf{f}}{\partial \mathbf{X}}\frac{\partial \mathbf{X}}{\partial \mathbf{X_0}}
\end{equation}

\begin{equation}
\drv \frac{\partial \mathbf{X}}{\partial \mathbf{X_0}}=
\frac{\partial \mathbf{f}}{\partial \mathbf{X}}
\frac{\partial \mathbf{X}}{\partial \mathbf{X_0}} \drv t
+
\frac{\partial \mathbf{G}}{\partial \mathbf{X}}
\frac{\partial \mathbf{X}}{\partial \mathbf{X_0}} \drv \mathbf{W_t}
\end{equation}

\begin{equation}
    \sigma_{t_0}^T(\mathbf{X})=\frac{1}{|T|}ln(\sqrt{\lambda_{max}(\Delta(\mathbf{X})}))
\end{equation}

Somewhere here mention random matrix theory \cite{potters_bouchaud_2020}.

%slide8 Stochastic Duffing oscillator
Another classical problem is the Duffing oscillator:
\begin{equation}
\ddot{x}=\alpha \dot{x}+\beta x+\gamma x^3 + \epsilon\cos{t}
\end{equation}
whose stochastic counterpart is: 
\begin{equation}
    \begin{cases}
    \drv X_t= \alpha Y_t \\
    \drv Y_t = (\beta X_t + \gamma X_t^3)\drv t +\epsilon X_t \drv \mathrm{W_t} 
\end{cases}
\end{equation}


\section{Results}

\begin{figure}[h]
    \centerline{\includegraphics[width=4cm]{juliagraphs.png}}
    \caption{This is example of the image in a column.}
        \label{fig:sample_figure}
    \end{figure}

\section{Conclusions and Recommendations}
Diffusion mechanisms arise both from chaotic behaviour and stochastic perturbations. But there is also an interplay.

Apply the methodology to systems in Celestial Mechanics: build on "Manzi, Boschmann. The Stochastic Three-body problem: Stochastic Resonances and Diffusion in small-body dynamics, Theory, models and simulations in Celestial Mechanics, 2022."

% **************GENERATED FILE, DO NOT EDIT**************

\bibliographystyle{juliacon}
\bibliography{ref.bib}


\end{document}

% Inspired by the International Journal of Computer Applications template
